\section{Metodología}\label{sec:metodologia-implementacio}

Para el desarrollo de la fase de implementación se utilizó un servidor tipo torre perteneciente a la infraestructura del grupo de investigación \ac{GRID}, sobre el cual se desplegó el entorno destinado a la ejecución de la arquitectura definida durante la etapa de diseño. Este enfoque permitió disponer de una infraestructura controlada para llevar a cabo la instalación, configuración, integración y validación de cada uno de los componentes del sistema.

Con el propósito de aprovechar los recursos disponibles y mantener el aislamiento de los servicios, se implementó un esquema de virtualización mediante la creación de las máquinas virtuales correspondientes a cada uno de los elementos de la solución. Asimismo, se realizó la asignación de los recursos de procesamiento, memoria y almacenamiento de acuerdo con los requisitos particulares de cada servicio, considerando aspectos relacionados con el rendimiento, la disponibilidad y la escalabilidad de la infraestructura.

De manera complementaria, se efectuó la configuración del esquema de red definido para el proyecto, incluyendo el direccionamiento de los equipos, la resolución de nombres, la definición de los nombres de dominio completamente calificados (\ac{FQDN}), la sincronización horaria y la aplicación de las reglas de seguridad necesarias para permitir el funcionamiento de los servicios implementados.

Finalmente, se llevó a cabo el proceso de instalación y configuración de la plataforma FreeIPA y de los servicios asociados a esta, verificando la correcta integración entre los distintos componentes y el cumplimiento de los requisitos funcionales y de seguridad establecidos para la infraestructura del grupo de investigación \ac{GRID}.
